\chapter{Fairest Edge Usage and Minimum Expected Overlap for Random Spanning Trees}
\label{chap:feu}

This chapter formalizes results from \feupaperref{}. It specializes the general
modulus theory of Chapter~\ref{chap:common} to spanning trees, and develops the
paper's two main optimization problems: minimum expected overlap (MEO) and
fairest edge usage (FEU).

\section{Random spanning trees, MEO, and FEU}

This section follows Sections 1.1--1.3 of the paper: it sets up random spanning
trees and states the two optimization problems, minimum expected overlap
(MEO) and fairest edge usage (FEU), whose relationship to spanning tree
modulus is the subject of the rest of the paper.

\begin{definition}[Law of a random spanning tree]
\label{def:spanning-tree-law}
\uses{def:spanning-tree}
Let $G = (V, E)$ be a finite, connected multigraph and let $\Gamma_G$ be its
(finite) set of spanning trees. A \emph{law} on $\Gamma_G$ is a probability
mass function $\mu \colon \Gamma_G \to [0,1]$ with $\sum_{\gamma \in \Gamma_G}
\mu(\gamma) = 1$; we write $\mu \in P(\Gamma_G)$. A random spanning tree
$\underline\gamma$ \emph{has law} $\mu$ (written $\underline\gamma \sim \mu$)
if $\mathbb P_\mu(\underline\gamma = \gamma) = \mu(\gamma)$ for every $\gamma
\in \Gamma_G$.
\end{definition}

\begin{definition}[Minimum expected overlap problem]
\label{def:meo}
\uses{def:spanning-tree-law}
For $\mu \in P(\Gamma_G)$, let $\underline\gamma, \underline\gamma' \sim \mu$
be independent. The \emph{expected overlap} of $\mu$ is
\[
  \mathbb E_\mu |\underline\gamma \cap \underline\gamma'|
  := \sum_{\gamma, \gamma' \in \Gamma_G} |\gamma \cap \gamma'| \, \mu(\gamma) \mu(\gamma').
\]
The \emph{minimum expected overlap (MEO) problem} is
\[
  \MEO(\Gamma_G) := \min_{\mu \in P(\Gamma_G)} \mathbb E_\mu |\underline\gamma \cap \underline\gamma'|.
\]
A law achieving this minimum is \emph{optimal} for $\MEO(\Gamma_G)$.
\end{definition}

\begin{definition}[Edge usage probability]
\label{def:edge-usage-probability}
\uses{def:spanning-tree-law}
For $\mu \in P(\Gamma_G)$ and $e \in E$, the \emph{edge usage probability} of
$e$ under $\mu$ is
\[
  \mathbb P_\mu(e \in \underline\gamma) := \sum_{\gamma \in \Gamma_G} \mathbf 1_{\{e \in \gamma\}} \, \mu(\gamma).
\]
\end{definition}

\begin{definition}[Fairest edge usage problem]
\label{def:feu}
\uses{def:edge-usage-probability}
The \emph{fairest edge usage (FEU) problem} asks for the edge usage
probability function $\eta \colon E \to [0,1]$ of minimal variance over $\mu
\in P(\Gamma_G)$:
\[
  \text{minimize } \operatorname{Var}(\eta)
  \quad \text{subject to} \quad
  \eta(e) = \mathbb P_\mu(e \in \underline\gamma) \ \forall e \in E, \quad \mu \in P(\Gamma_G).
\]
We write $\FEU(\Gamma_G)$ for this problem.
\end{definition}

\begin{definition}[Fair and forbidden trees]
\label{def:fair-tree}
\uses{def:meo}
A spanning tree $\gamma \in \Gamma_G$ is a \emph{fair tree} if $\mu^*(\gamma) >
0$ for some law $\mu^*$ optimal for $\MEO(\Gamma_G)$. The set of fair trees is
denoted $\Gamma_G^f$. A tree $\gamma \in \Gamma_G \setminus \Gamma_G^f$ (if any
exist) is a \emph{forbidden tree}.
\end{definition}

\section{Spanning tree modulus and the MEO problem}
\label{sec:spanning-tree-modulus}

This section follows Section 1.7 of the paper and connects the general
modulus framework of Chapter~\ref{chap:common} back to spanning trees,
recovering MEO and FEU as modulus problems. From here on the paper restricts
to $p = 2$ and unweighted graphs ($\sigma \equiv 1$), and we do the same.

\begin{definition}[Feasible partition]
\label{def:feasible-partition}
\uses{def:multigraph}
A \emph{feasible partition} of $G = (V, E)$ is a partition $P = \{V_1, \dots,
V_{k_P}\}$ of $V$ into $k_P \ge 2$ parts such that each induced subgraph
$G(V_i)$ is connected. Its edge set $E_P \subseteq E$ consists of the edges
joining vertices belonging to different parts.
\end{definition}

\begin{theorem}[Fulkerson dual of spanning trees]
\label{thm:spanning-tree-fulkerson-dual}
\uses{def:feasible-partition, def:fulkerson-dual, def:spanning-tree}
Let $\Gamma_G$ be the family of spanning trees of $G = (V, E)$, with usage
given by the indicator vectors $\mathbf 1_{\{e \in \gamma\}}$. The Fulkerson
dual family $\widehat{\Gamma_G}$ is the set of vectors $\frac{1}{k_P - 1}
\mathbf 1_{E_P}$, ranging over all feasible partitions $P$ of $G$.
\end{theorem}

\begin{theorem}[Spanning tree modulus, MEO, and FEU]
\label{thm:modulus-meo-feu}
\uses{def:modulus, def:meo, def:feu, thm:spanning-tree-fulkerson-dual}
Let $\Gamma = \Gamma_G$ be the family of spanning trees of $G = (V, E)$ and let
$\widehat\Gamma$ be its Fulkerson dual. Densities $\rho, \eta \in \mathbb
R_{\ge 0}^E$ and a law $\mu \in P(\Gamma)$ are simultaneously optimal for
$\Mod_2(\Gamma)$, $\Mod_2(\widehat\Gamma)$, and $\MEO(\Gamma)$ respectively if
and only if
\[
  \rho \in \Adm(\Gamma), \quad \eta = \usageN^T \mu, \quad
  \eta(e) = \frac{\rho(e)}{\Mod_2(\Gamma)} \ \forall e \in E, \quad
  \mu(\gamma) (1 - \ell_\rho(\gamma)) = 0 \ \forall \gamma \in \Gamma.
\]
In particular,
\[
  \MEO(\Gamma) = \Mod_2(\widehat\Gamma) = \Mod_2(\Gamma)^{-1}.
\]
\end{theorem}

\begin{corollary}[Weak duality for modulus]
\label{cor:weak-duality}
\uses{def:modulus}
Let $\Gamma$ be a finite family of objects on $G$ with Fulkerson dual
$\widehat\Gamma$. For any $\rho \in \Adm(\Gamma)$ and $\eta \in
\Adm(\widehat\Gamma)$,
\[
  E_2(\rho) E_2(\eta) \ge 1, \qquad \text{where } E_2(\rho) := \sum_{e \in E} \sigma(e) \rho(e)^2,
\]
with equality if and only if $\rho$ and $\eta$ are optimal for
$\Mod_2(\Gamma)$ and $\Mod_2(\widehat\Gamma)$ respectively.
\end{corollary}
